\begin{textsample}{POS Dim 6 – vem\_ed\_30 – Score 10.00 – vem\_ed\_30/t158.txt}  \label{ex:f6_pos_100}
PCI aborda \textbf{emissões} de \textbf{poluentes} em \textbf{veículos} No primeiro semestre de 2025, o professor Marco Aurélio Fróes, do curso de Mecânica Automobilística da Fatec Santo André, desenvolveu um Projeto Colaborativo Internacional (PCI) com o professor Javier Pavón, da DUOC UC (Chile).
O projeto permitiu aos estudantes brasileiros e \textbf{chilenos} conhecerem \textbf{normas} internacionais utilizadas na certificação de \textbf{emissões} de \textbf{poluentes} em \textbf{veículos} leves, praticarem \textbf{ensaios} em laboratório e compararem os resultados.
O professor Orlando de Salvo Júnior, especialista em \textbf{emissões} \textbf{veiculares}, foi convidado pelo professor Fróes para auxiliar no entendimento dos conteúdos técnicos.
Esse foi o primeiro PCI de Fróes. “Fiquei um pouco perdido inicialmente, pois me faltava experiência, mas a professora Regiane orientou e conduziu muito bem, trazendo segurança”, comenta, referindo-se a Regiane Souza Camargo Moreira, responsável pela coordenação dos PCIs em espanhol na área de Apoio à Internacionalização do Ensino Superior da Divisão de Extensão e Pesquisa (Depes) da Coordenadoria Geral de Ensino Superior de Graduação (CGESG) do CPS.
As interações dos grupos de alunos se deram por WhatsApp, videochamadas no Zoom ou Google Meet, e os estudos desenvolvidos foram compartilhados no \textbf{Padlet} (em \textbf{português} e espanhol para driblar as barreiras idiomáticas).
No entender de Fróes, “a troca de experiências e de realidades com uma instituição de ensino internacional foi o principal ganho.
Foi possível ver, na prática, a aplicação de uma \textbf{norma} internacional para \textbf{homologação} \textbf{veicular} com todos os critérios de segurança e condicionamento, bem como estabelecer comparações com os procedimentos de \textbf{ensaio} de potência e \textbf{torque}”. continuação O professor Javier Pavón, em sua reflexão final sobre o projeto publicada no \textbf{Padlet}, avaliou: “Durante todo o processo, os alunos adquiriram novos conhecimentos sobre \textbf{normas} e procedimentos relacionados a testes de motores em dinamômetro, o que lhes permitiu ampliar seus conhecimentos técnicos e aplicar metodologias de teste de acordo com as \textbf{normas} internacionais.
Além disso, o trabalho em conjunto com os colegas brasileiros ofereceu a possibilidade de conhecer diferentes culturas e experiências no setor automotivo, promovendo uma visão global sobre a \textbf{homologação} e certificação de \textbf{veículos}”.

% matched lemmas: chileno, emissão, ensaio, homologação, norma, padlet, poluente, português, torque, veicular, veículo
\end{textsample}
